\documentclass[12pt]{article}
\usepackage[utf8]{inputenc}
\usepackage{amsmath}
\usepackage{amsfonts}
\usepackage{amssymb}
\usepackage{amsthm}
\usepackage{geometry}
\usepackage{parskip}
\usepackage{xcolor}
\usepackage{graphicx}
\usepackage{enumitem}
\usepackage{hyperref}

\geometry{a4paper, margin=1in}

% Custom command for a "Key Insight" box
\newcommand{\insight}[1]{%
    \begin{center}
    \fcolorbox{black}{gray!10}{\begin{minipage}{0.92\textwidth}
    \textbf{Key Insight:} #1
    \end{minipage}}
    \end{center}
}

\graphicspath{ {./images/} }

\title{KCPC Workshop 2 \\ Solutions: Techniques for Investigation}
\author{King's Competitive Programming Club}
\date{}

\begin{document}

\maketitle

\noindent
These official solutions accompany Workshop 2 of the KCPC 2nd Semester series. They illustrate fundamental mathematical and algorithmic heuristics: spatial reasoning, functional recurrences, induction, visual invariants, change of reference frame, symmetry, the extreme principle, and competitive programming game analysis.

\vspace{1em}
\hrule
\vspace{1em}

\section{Psychological Strategies}

\subsection{Pat's Dilemma}

\insight{Do not artificially constrain your thinking to two dimensions (Flatland). An object whose length exceeds 1 metre fits into a box whose individual dimensions do not exceed 1 metre by placing it along the space diagonal.}

\textbf{Solution:}
Pat should place the sword diagonally inside a rectangular cubic box measuring $1\text{ m} \times 1\text{ m} \times 1\text{ m}$.

Every linear dimension of this box (length = 1 m, width = 1 m, height = 1 m) satisfies the baggage regulation of not exceeding 1 metre. By the three-dimensional Pythagorean theorem, the interior space diagonal of this box is:
\[
d = \sqrt{1^2 + 1^2 + 1^2} = \sqrt{3} \approx 1.732\text{ metres}.
\]
Because $1.732\text{ m} > 1.5\text{ m}$, the 1.5-metre sword fits comfortably along the diagonal without bending or disassembling.

\subsection{Thinking in 3D}

\insight{When two polyhedra are glued together, check whether adjacent faces on either side of the joint become coplanar. If the dihedral angles sum to $180^\circ$, adjacent triangles merge into a single flat quadrilateral.}

\textbf{Answer:} Exactly 5 faces.

\textbf{Explanation:}
The square pyramid initially has 5 faces (1 square base and 4 equilateral triangular lateral faces). The regular tetrahedron has 4 equilateral triangular faces. All edge lengths are 1.

When the tetrahedron is glued to the pyramid along one triangular face, that interface face is buried in the interior, removing 1 face from each solid:
\[
5 + 4 - 2 = 7 \text{ faces initially.}
\]
However, we must analyze the dihedral angle between the faces meeting at the boundaries of the glued triangle:
\begin{itemize}
    \item In a regular tetrahedron of unit edge length, the angle between any two adjacent triangular faces is $\theta_{\text{tet}} = \arccos(1/3) \approx 70.53^\circ$.
    \item In a square pyramid with equilateral triangular sides of unit length, the dihedral angle between two adjacent triangular faces sharing a slant edge is also $\theta_{\text{pyr}} = \arccos(-1/3) = 180^\circ - \arccos(1/3) \approx 109.47^\circ$.
    \item At the two edges adjacent to the glued face, the angle between the lateral face of the pyramid and the adjacent face of the tetrahedron is:
    \[
    \theta_{\text{tet}} + \theta_{\text{pyr-edge}} = \arccos(1/3) + (180^\circ - \arccos(1/3)) = 180^\circ.
    \]
\end{itemize}
Because this combined angle is exactly $180^\circ$, each of the two adjacent tetrahedron faces lies in the exact same plane as its adjoining pyramid face. The two pairs of coplanar triangles merge into two flat rhombi (quadrilaterals).

This reduces the face count by 2 more:
\[
7 - 2 = 5 \text{ faces.}
\]
The resulting solid has 1 square face, 2 equilateral triangular faces, and 2 rhombus faces, for a total of 5 faces.

\begin{center}
\includegraphics[width=0.55\linewidth]{polyhedra.png}
\end{center}

\section{Wishful Thinking and Make It Easier}

\subsection{Long Sequence}

\insight{When given an intimidating multi-variable functional equation, evaluate the smallest possible indices first ($m=n=0$, then $n=0$) to uncover the underlying recurrence, formulate a conjecture, and complete the proof by strong mathematical induction.}

\textbf{Answer:} $a_{6767} = 6767^2 = 45{,}792{,}289$.

\textbf{Complete Proof:}
The given relation is:
\begin{equation*}
    a_{m+n} + a_{m-n} = \frac{1}{2}(a_{2m} + a_{2n}) \quad \text{for all } m \ge n \ge 0.
\end{equation*}

\textbf{Step 1: Discovering the functional structure.}
\begin{itemize}
    \item Substitute $m = 0, n = 0$:
    \[
    a_0 + a_0 = \frac{1}{2}(a_0 + a_0) \implies 2a_0 = a_0 \implies a_0 = 0.
    \]
    \item Substitute $n = 0$ for any $m \ge 0$:
    \[
    a_m + a_m = \frac{1}{2}(a_{2m} + a_0) = \frac{1}{2}a_{2m} \implies a_{2m} = 4a_m.
    \]
    \item Now replace $a_{2m}$ with $4a_m$ and $a_{2n}$ with $4a_n$ in the original equation:
    \[
    a_{m+n} + a_{m-n} = \frac{1}{2}(4a_m + 4a_n) = 2(a_m + a_n).
    \]
    This is the classical parallelogram law functional equation (Jordan--von Neumann).
\end{itemize}

\textbf{Step 2: Calculating small values.}
\begin{itemize}
    \item $a_0 = 0 = 0^2$.
    \item $a_1 = 1 = 1^2$ (given).
    \item $a_2 = a_{2(1)} = 4a_1 = 4 = 2^2$.
    \item Setting $m = 2, n = 1$:
    \[
    a_3 + a_1 = 2(a_2 + a_1) \implies a_3 + 1 = 2(4 + 1) = 10 \implies a_3 = 9 = 3^2.
    \]
\end{itemize}

\textbf{Step 3: Rigorous inductive proof.}
We prove by strong induction that $a_k = k^2$ for all non-negative integers $k$.

\textit{Base Cases:} The identity holds for $k = 0, 1, 2, 3$.

\textit{Inductive Step:} Assume $a_j = j^2$ holds for all integers $0 \le j \le k$ where $k \ge 1$. We compute $a_{k+1}$ by setting $m = k$ and $n = 1$:
\[
a_{k+1} + a_{k-1} = 2(a_k + a_1).
\]
By the inductive hypothesis, $a_{k-1} = (k-1)^2$, $a_k = k^2$, and $a_1 = 1$. Substituting these in:
\[
a_{k+1} + (k-1)^2 = 2(k^2 + 1)
\]
\[
a_{k+1} = 2k^2 + 2 - (k^2 - 2k + 1) = k^2 + 2k + 1 = (k+1)^2.
\]
By mathematical induction, $a_n = n^2$ for all integers $n \ge 0$.

Therefore, for $n = 6767$:
\[
a_{6767} = 6767^2 = 45{,}792{,}289.
\]

\subsection{An \textit{Algorithm?}}

\insight{The recurrence $f(2n) = f(n)$ and $f(2n+1) = f(n)+1$ defines the population count (popcount) of the binary representation of $n$.}

\textbf{Solution:}
$f(n)$ equals the number of ones in the binary representation of $n$ (the Hamming weight or popcount of $n$).

\textbf{Explanation:}
Consider $n$ written in base 2:
\begin{itemize}
    \item The base case is $f(1) = 1$, and $1 = (1)_2$ has exactly one $1$.
    \item Even numbers $2n$ end with a $0$ in binary: $(2n)_2 = (n)_2 \, 0$. Shifting the binary string to the right by one position removes the trailing $0$, leaving the number of ones unchanged. Thus, $f(2n) = f(n)$.
    \item Odd numbers $2n+1$ end with a $1$ in binary: $(2n+1)_2 = (n)_2 \, 1$. Truncating the trailing $1$ leaves $(n)_2$, so the count of ones increases by 1: $f(2n+1) = f(2n) + 1 = f(n) + 1$.
\end{itemize}
By induction on the number of bits, $f(n)$ counts the set bits in $n$'s binary expansion.

\section{Induction}

\insight{To prove an inequality by induction, identify the multiplicative step from the base hypothesis to the goal: $(n+1)! = (n+1) \cdot n! > (n+1) \cdot 2^n > 2 \cdot 2^n = 2^{n+1}$.}

\textbf{Problem:} Prove that $n! > 2^n$ for all integers $n \ge 4$.

\begin{proof}
We proceed by mathematical induction on $n \ge 4$.

\textbf{Base Case ($n = 4$):}
For $n = 4$, the left-hand side is $4! = 24$, and the right-hand side is $2^4 = 16$. Since $24 > 16$, the base case holds.

\textbf{Inductive Hypothesis:}
Assume that for some integer $k \ge 4$, the inequality holds:
\[
k! > 2^k.
\]

\textbf{Inductive Step:}
We wish to prove that $(k+1)! > 2^{k+1}$.
Multiplying both sides of the inductive hypothesis by $k + 1 > 0$ preserves the inequality:
\[
(k+1)! = (k+1) \cdot k! > (k+1) \cdot 2^k.
\]
Because $k \ge 4$, we have $k + 1 \ge 5 > 2$. Therefore:
\[
(k+1) \cdot 2^k > 2 \cdot 2^k = 2^{k+1}.
\]
Combining these inequalities gives:
\[
(k+1)! > 2^{k+1}.
\]
This completes the inductive step. By the principle of mathematical induction, $n! > 2^n$ for all integers $n \ge 4$.
\end{proof}

\section{Draw a Picture}

\insight{A square of side length $2 \sum_{i=1}^n i = n(n+1)$ can be decomposed into $n$ concentric gnomon rings, where the $i$-th ring consists of $4i$ small squares of side length $i$, having total area $4i^3$.}

\textbf{Visual Proof:}
Consider a large square whose side length is twice the $n$-th triangular number:
\[
S = 2(1 + 2 + 3 + \dots + n) = 2 \cdot \frac{n(n+1)}{2} = n(n+1).
\]
The total area of this square is:
\[
\text{Area} = S^2 = [n(n+1)]^2 = 4 \left( \sum_{i=1}^n i \right)^2.
\]

\begin{center}
\includegraphics[width=0.55\linewidth]{cubes.png}
\end{center}

We now partition this large square into $n$ concentric square rings (gnomons) centered at the origin:
\begin{itemize}
    \item The 1st level (center) consists of 4 squares of side length 1, with total area $4 \times 1 \times 1^2 = 4(1^3)$.
    \item The 2nd level adds a ring around the center consisting of $4 \times 2 = 8$ squares of side length 2, with total area $4 \times 2 \times 2^2 = 4(2^3)$.
    \item In general, the $i$-th concentric layer consists of $4i$ squares of side length $i$. Its area is:
    \[
    \text{Area}_i = 4i \times i^2 = 4i^3.
    \]
\end{itemize}
Summing the areas of all $n$ layers yields:
\[
\text{Area} = \sum_{i=1}^n 4i^3 = 4 \sum_{i=1}^n i^3.
\]
Equating the two expressions for the area:
\[
4 \sum_{i=1}^n i^3 = 4 \left( \sum_{i=1}^n i \right)^2 \implies \sum_{i=1}^n i^3 = \left( \sum_{i=1}^n i \right)^2 = \frac{n^2(n+1)^2}{4}.
\]

\section{Recast}

\insight{Coloring is an invariant tool. A standard domino always covers one black square and one white square. If a grid's black and white counts differ, it cannot be tiled by dominoes.}

\textbf{Answer:} No, it is impossible to tile the mutilated board.

\textbf{Proof:}
Color the $8 \times 8$ chessboard with alternating black and white squares in the standard fashion (32 black squares and 32 white squares).

Notice that diagonally opposite corner squares on a chessboard always share the same color. Without loss of generality, suppose the two removed corners are both white.
The remaining 62-square mutilated board contains:
\begin{itemize}
    \item 32 black squares,
    \item 30 white squares.
\end{itemize}
Every $2 \times 1$ domino placed on the board covers exactly two adjacent squares. Because horizontally or vertically adjacent squares always have opposite colors, every placed domino covers exactly \textbf{one black square and one white square}.

Therefore, any valid tiling using 31 non-overlapping dominoes must cover exactly 31 black squares and 31 white squares. Because the mutilated board has 32 black and only 30 white squares, no such tiling can possibly exist.

\section{Change Your Point of View}

\insight{Instead of setting up algebraic equations, shift into the rotating reference frame of the hour hand. In 12 hours, the minute hand completes 11 relative laps around the clock face.}

\textbf{Answer:} At exactly $1:05\frac{5}{11}$ (or $1\text{ hour, } 5\text{ minutes, and } 27\frac{3}{11}\text{ seconds}$).

\textbf{Explanation:}
Consider the relative angular speeds of the two hands:
\begin{itemize}
    \item The minute hand travels at 1 revolution per hour ($360^\circ/\text{hr}$).
    \item The hour hand travels at $1/12$ revolution per hour ($30^\circ/\text{hr}$).
    \item In the reference frame of the hour hand, the minute hand moves at a relative speed of:
    \[
    v_{\text{rel}} = 1 - \frac{1}{12} = \frac{11}{12} \text{ revolutions per hour}.
    \]
\end{itemize}
During the 12 hours between 12:00 noon and 12:00 midnight, the minute hand makes 12 laps while the hour hand makes 1 lap. Thus, the minute hand laps (overtakes) the hour hand exactly:
\[
12 - 1 = 11 \text{ times}.
\]
Because both hands travel at constant angular speeds, these 11 meeting points are distributed with perfect rotational symmetry around the clock at equal time intervals:
\[
\Delta t = \frac{12 \text{ hours}}{11} = 1\frac{1}{11} \text{ hours}.
\]
The first meeting after 12 o'clock occurs at:
\[
t = 1\frac{1}{11} \text{ hours} = 1 \text{ hour} + \frac{60}{11} \text{ minutes} = 1 \text{ hour} + 5\frac{5}{11} \text{ minutes} \approx 1:05:27.27.
\]

\section{Symmetry}

\subsection{Grandma's Cabin}

\insight{Heron's shortest-path reflection principle: reflecting one point across the boundary converts a broken path into a single straight line.}

\textbf{Answer:} Exactly 13 miles.

\textbf{Explanation:}
Let the east-west stream be the $x$-axis ($y = 0$).
\begin{itemize}
    \item Your cabin is at $Y = (0, 2)$.
    \item Grandma's cabin is 12 miles west and 1 mile north of your cabin, so $G = (-12, 3)$.
\end{itemize}
We must choose a point $A = (x, 0)$ on the stream to minimize the total walking distance $|YA| + |AG|$.

\begin{center}
\includegraphics[width=0.45\linewidth]{grandma.png}
\end{center}

Reflect Grandma's cabin across the stream ($y = 0$) to point $G' = (-12, -3)$.
For any point $A$ on the stream, by symmetry across the $x$-axis, $|AG| = |AG'|$.
Therefore, the path distance is:
\[
|YA| + |AG| = |YA| + |AG'|.
\]
By the triangle inequality in Euclidean space, the shortest path between $Y$ and $G'$ is the direct straight line segment $YG'$.
The horizontal separation between $Y(0, 2)$ and $G'(-12, -3)$ is $\Delta x = 12$ miles, and the vertical separation is $\Delta y = 2 - (-3) = 5$ miles.
Applying the Pythagorean theorem:
\[
\text{Minimum Distance} = \sqrt{(\Delta x)^2 + (\Delta y)^2} = \sqrt{12^2 + 5^2} = \sqrt{144 + 25} = \sqrt{169} = 13 \text{ miles}.
\]

\subsection{Geometric Symmetry}

\insight{Rotate the inner square by $45^\circ$ so that its four vertices touch the points of tangency between the circle and the outer square.}

\textbf{Answer:} The ratio of the area of the smaller square to the larger square is $1:2$ (or $1/2$).

\begin{center}
\includegraphics[width=0.28\linewidth]{symm.png}
\end{center}

\textbf{Explanation:}
Rotate the smaller inscribed square by $45^\circ$ inside the circle:
\begin{itemize}
    \item Its four vertices now touch the midpoints of the four edges of the outer square (the points where the circle touches the outer square).
    \item The four corners of the outer square outside the rotated inner square form four right-angled triangles whose total area equals the area of the inner square itself.
    \item Hence, the outer square is decomposed into 8 identical right triangles, 4 of which constitute the inner square.
\end{itemize}
Thus, the inner square has exactly half the area of the outer square:
\[
\frac{\text{Area}(\text{Inner})}{\text{Area}(\text{Outer})} = \frac{1}{2}.
\]

\subsection{Symmetry in Probability}

\insight{By symmetry, $k$ randomly chosen division points partition a sequence or interval into $k+1$ statistically identical segments. Linearity of expectation applies without requiring complex conditional probabilities.}

\textbf{Answers:}
\begin{enumerate}[label=(\alph*)]
    \item \textbf{Expected cards until the first ace:} Exactly $10.6$ cards (or $\frac{53}{5}$).

    \textbf{Derivation:}
    The 4 aces in the 52-card deck partition the 48 non-ace cards into 5 intervals:
    \begin{enumerate}[label=\arabic*.]
        \item Cards before the 1st ace,
        \item Cards between the 1st and 2nd aces,
        \item Cards between the 2nd and 3rd aces,
        \item Cards between the 3rd and 4th aces,
        \item Cards after the 4th ace.
    \end{enumerate}
    By permutation symmetry, each of the 48 non-ace cards is equally likely to land in any of the 5 intervals.
    Therefore, by linearity of expectation, the expected number of non-ace cards in each interval is:
    \[
    \mathbb{E}[\text{non-aces before 1st ace}] = \frac{48}{5} = 9.6.
    \]
    The total cards dealt up to and including the first ace is the number of non-aces preceding it plus the first ace itself:
    \[
    \mathbb{E}[\text{cards dealt}] = 9.6 + 1 = 10.6 = \frac{53}{5}.
    \]

    \item \textbf{Putnam 1992 (Four random points on a sphere):} The probability is $\dfrac{1}{8}$.

    \textbf{Derivation:}
    Let the 4 points chosen independently and uniformly at random on the unit sphere be $P_1, P_2, P_3, P_4$.
    Almost surely, no four points are coplanar, and no two points are antipodal.

    For each point $P_i$, consider its antipodal point $-P_i$. The 4 pairs of antipodal points $\{P_i, -P_i\}$ define 4 lines through the origin, which generate $2^4 = 16$ possible tetrahedra formed by selecting one vertex from each pair.

    The centre of the sphere (the origin) is contained inside the convex hull of 4 points if and only if the origin cannot be separated from the points by a plane through the origin. The 4 great circles (planes through the origin perpendicular to the lines) partition the sphere into regions. By projective duality and Euler's formula for spherical arrangements, exactly 2 of the 16 tetrahedra contain the origin in their interior.

    Because all 16 sign assignments $(\pm P_1, \pm P_2, \pm P_3, \pm P_4)$ are equally probable by spherical symmetry, the probability is:
    \[
    \mathbb{P}(\text{origin inside tetrahedron}) = \frac{2}{16} = \frac{1}{8}.
    \]
\end{enumerate}

\section{The Extreme Principle}

\insight{To prove a global geometric property, assume the contrary and inspect an extremal element (minimal distance, minimal non-zero area, or minimal distance to a line). Constructing a strictly smaller non-zero element yields an immediate contradiction.}

\textbf{Problem:} Prove that all points in $B \cup W$ lie on a single straight line.

\begin{proof}
Suppose, for the sake of contradiction, that the points in $B \cup W$ do not all lie on a single line.

Then there exists at least one pair of points of the same color, say $P, Q \in B$ (without loss of generality, black), and a point $R \in B \cup W$ not lying on the line $L_{PQ}$ connecting $P$ and $Q$.
Thus, the distance from $R$ to the line $L_{PQ}$ is strictly positive: $d(R, L_{PQ}) > 0$.

Because $B \cup W$ is a finite set, there exist only finitely many triples $(P, Q, R)$ where $P, Q$ have the same colour and $R \notin L_{PQ}$.
By the well-ordering principle, there must exist a configuration $(P_0, Q_0, R_0)$ that strictly \textbf{minimizes} the positive perpendicular distance:
\[
h_{\min} = d(R_0, L_{P_0 Q_0}) > 0.
\]
By the given problem condition, the segment $P_0 Q_0$ connects two black points, so it must contain at least one point of the other colour, say a white point $W_1 \in P_0 Q_0$.

Now examine the color of $R_0$:
\begin{itemize}
    \item \textbf{Case 1: $R_0$ is white.}
    Then $R_0$ and $W_1$ are two points of the same color (white). The line segment $R_0 W_1$ must contain a point of the other color (black), say $B_1 \in R_0 W_1$ with $B_1 \neq R_0$ and $B_1 \neq W_1$.
    Consider the distance from $B_1$ to the line $L_{P_0 Q_0}$. Because $B_1$ lies strictly in the interior of the line segment $R_0 W_1$, its height above $L_{P_0 Q_0}$ is strictly between $0$ and $h_{\min}$:
    \[
    0 < d(B_1, L_{P_0 Q_0}) < d(R_0, L_{P_0 Q_0}) = h_{\min}.
    \]
    The pair $P_0, Q_0$ are both black, and $B_1 \notin L_{P_0 Q_0}$ is a point whose distance to $L_{P_0 Q_0}$ is strictly smaller than $h_{\min}$. This directly contradicts the minimality of $h_{\min}$!

    \item \textbf{Case 2: $R_0$ is black.}
    Then $R_0$ and $P_0$ (or $R_0$ and $Q_0$) are two points of the same colour (black). Consider the line $L_{R_0 P_0}$. The projection of $W_1$ onto this line yields a strictly smaller positive perpendicular distance to the line than $h_{\min}$ (since $W_1$ lies strictly between $P_0$ and $Q_0$). This again contradicts the minimality of $h_{\min}$.
\end{itemize}
In all cases, assuming the points do not lie on a single line leads to a contradiction. Therefore, all points in $B \cup W$ must lie on a single straight line.
\end{proof}

\section{Advanced Homework (Competitive Programming)}

\subsection{ARC 122, Problem C --- Calculator}

\insight{Operations 3 and 4 alternate to generate Fibonacci numbers. Using Zeckendorf's theorem or reverse Euclidean division, any positive integer $N \le 10^{18}$ can be decomposed into at most $\approx 86$ Fibonacci steps, requiring $\le 130$ total operations.}

\textbf{Solution Analysis:}
We are given initial values $x = 0, y = 0$ and want to reach $x = N$ in at most 130 operations:
\begin{itemize}
    \item Op 1: $x \leftarrow x + 1$
    \item Op 2: $y \leftarrow y + 1$
    \item Op 3: $x \leftarrow x + y$
    \item Op 4: $y \leftarrow x + y$
\end{itemize}
Notice that if we alternate Op 3 and Op 4:
\[
(1, 0) \xrightarrow{\text{Op 4}} (1, 1) \xrightarrow{\text{Op 3}} (2, 1) \xrightarrow{\text{Op 4}} (2, 3) \xrightarrow{\text{Op 3}} (5, 3) \xrightarrow{\text{Op 4}} (5, 8) \dots
\]
The values generated are precisely the consecutive Fibonacci numbers $F_k$.
The 86th Fibonacci number is $F_{86} \approx 2.8 \times 10^{18} > 10^{18}$, which comfortably exceeds the maximum possible value of $N$.

\textbf{Algorithmic Construction:}
We work backwards using the greedy Fibonacci representation (Zeckendorf's theorem):
\begin{enumerate}
    \item Precompute Fibonacci numbers up to $F_K$ such that $F_K \le N < F_{K+1}$, where $K \le 86$.
    \item Express $N$ greedily as a sum of non-consecutive Fibonacci numbers:
    \[
    N = \sum_{j \in S} F_j, \quad S \subseteq \{1, 2, \dots, K\}.
    \]
    \item Simulate the process forward from step 1 to $K$:
    \begin{itemize}
        \item At step $j$, if $j$ is odd, we execute Op 3 ($x \leftarrow x + y$). If $j$ is even, we execute Op 4 ($y \leftarrow x + y$).
        \item If $j \in S$, we insert an increment (Op 1 if $j$ is odd, Op 2 if $j$ is even) to add the required unit corresponding to $F_j$.
    \end{itemize}
    \item Because no two consecutive Fibonacci numbers are selected in the Zeckendorf representation, the number of increments is at most $\lceil 86 / 2 \rceil = 43$.
\end{enumerate}
The total number of operations performed is at most $86 + 43 = 129 \le 130$, easily fitting within the problem's strict constraint.

\subsection{ARC 131, Problem C --- Zero XOR}

\insight{Let $S = A_1 \oplus \dots \oplus A_N$. If $A_i = S$ for some $i$, Player 1 wins immediately on move 1. Otherwise, analyze the parity: when $N$ is odd, the bitwise sum of all possible remaining XOR states is invariant, guaranteeing a winning strategy for Player 1.}

\textbf{Complete Strategy \& Answer:}
Compute the total bitwise XOR sum:
\[
S = \bigoplus_{i=1}^N A_i = A_1 \oplus A_2 \oplus \dots \oplus A_N.
\]
A player wins immediately if, after eating cookie $A_i$, the XOR sum of the remaining $N-1$ cookies is 0:
\[
S \oplus A_i = 0 \iff A_i = S.
\]

\textbf{Case Analysis:}
\begin{enumerate}
    \item \textbf{Direct Win on Turn 1:}
    If there exists an index $i \in \{1, \dots, N\}$ such that $A_i = S$, Player 1 chooses cookie $A_i$ on turn 1. The remaining XOR sum is 0, so Player 1 wins immediately.

    \item \textbf{General Case (No $A_i = S$):}
    \begin{itemize}
        \item \textbf{When $N$ is odd:}
        Player 1 is guaranteed to win!
        Consider the sum of all $N$ possible XOR states after one cookie is removed:
        \[
        \bigoplus_{i=1}^N (S \oplus A_i) = \left( \bigoplus_{i=1}^N S \right) \oplus \left( \bigoplus_{i=1}^N A_i \right) = (N \bmod 2) \cdot S \oplus S = S \oplus S = 0.
        \]
        Because the total XOR of all $N$ available moves is 0, Player 2 cannot force a loss for Player 1. If every move by Player 1 handed a win to Player 2, all $N$ states would have to be losing states for Player 2, which contradicts parity. As long as $N$ is odd, Player 1 can always maintain an optimal path to victory.

        \item \textbf{When $N$ is even (and no $A_i = S$):}
        Player 2 (Matthew) has the winning strategy. Any cookie Player 1 eats leaves an odd number of cookies $N - 1$, from which Player 2 assumes the winning role of the odd-parity player.
    \end{itemize}
\end{enumerate}
\textbf{Final Algorithm:}
Player 1 wins if and only if $N$ is odd OR there exists some $i \in \{1, \dots, N\}$ such that $A_i = S$. Otherwise, Matthew (Player 2) wins. The solution runs in $O(N)$ time and $O(1)$ auxiliary space.

\vspace{1em}
\hrule
\vspace{1em}

\section*{Attribution \& References}
\begin{itemize}
    \item Puzzles in Sections 1, 2, 6, 7, and 8 are adapted from classical problem-solving literature, primarily Paul Zeitz's \textit{The Art and Craft of Problem Solving} (John Wiley \& Sons) and Martin Gardner's \textit{Mathematical Games} columns.
    \item The Mutilated Chessboard problem is due to Max Black (1946) and popularized by Solomon W. Golomb and Martin Gardner.
    \item The 4-points-on-a-sphere problem appeared on the 53rd William Lowell Putnam Mathematical Competition (1992, Problem A6).
    \item Problems in Section 9 are official AtCoder Regular Contest tasks (ARC 122 C by editorialist \textit{beet}, ARC 131 C by editorialist \textit{maroonrk}).
\end{itemize}

\end{document}
